\section{Coherent descent and Euler characteristics}\label{sec:descent}

The matrices in Section~\ref{sec:curves} are defined over a function
field and need not themselves generate a division algebra.  We now
combine them with the parameter complete intersection to obtain a
division algebra that also contains the previous coefficient algebra.
The main step in this section converts the reduced rank of an arbitrary
module into an Euler characteristic.  Section~\ref{sec:index} uses that
conversion to prove the required rank divisibility.

\begin{theorem}\label{thm:division}
Let \(k\) be a finitely generated extension of the original field
\(\C\), and let \(\Delta_0\) be a central division algebra over \(k\)
of degree \(e_0\).  Fix an integer \(b\geq2\), a \(k\)-vector space
\(U\) of dimension at least three, and a subspace
\[
W\subseteq\Sym^b(U^\vee),\qquad s=\dim_k W,
\]
such that every nonzero member of \(W_{\bar k}\) has an irreducible
factor of multiplicity one.  Choose integers \(h,M\) in this order,
with
\begin{equation}\label{eq:division-parameters}
h\geq
\max\bigl(\{2\}\cup
\{\,v_p(e_0)s(p^{v_p(b)}-1):p\mid e_0\,\}\bigr),
\qquad M\geq2h+4s+10.
\end{equation}
Here \(v_p\) is the \(p\)-adic valuation, and the prime-indexed set is
empty if \(e_0=1\).
Define
\[
Y=\left\{(\ell_1,\ldots,\ell_M)\in U^{\oplus M}:
          \sum_{m=1}^M w(\ell_m)=0\quad(w\in W)\right\},
\qquad X=\PP(Y).
\]

Take the curve \(C\), group \(G\), and linearized line bundle \(L\)
of Proposition~\ref{prop:binary-curves} for this \(b\).  Write
\(e=|G|=b^b\) and \(g_T=b-1\), and put
\[
B=\prod_{m=2}^M\Pic^0(C),\qquad
H=G^{M-1},\qquad e_1=|H|.
\]
All products and curves here are over \(k\), and \(H\) acts on the
corresponding Picard factors by transport from the curve.
Choose arbitrary characters
\(\rho_1,\ldots,\rho_M:H\to\mu_b\), and let \(H\) act on the
points of the \(m\)-th block of \(Y\) by \(\rho_m\).
Let \(R\) be an affine space with a faithful linear \(H\)-action.
Set
\[
K=k(B\times X\times R),\qquad F=K^H.
\]
For the generic tuple of degree-zero line bundles \(A_m\), let
\[
E=\bigotimes_{m=2}^M H^0(C_K,A_m\otimes L^{g_T}).
\]
Transport induces a semilinear \(H\)-action on \(\End_K(E)\).
Its fixed algebra
\(\Delta_1=\End_K(E)^H\) is central simple of degree \(e_1\)
over \(F\), and
\begin{equation}\label{eq:division-algebra}
\Delta=(\Delta_0\otimes_k F)\otimes_F\Delta_1
\end{equation}
is a central division algebra of degree \(e_0e_1\).
Its scalar extension to \(K\) identifies with
\(\Delta_0\otimes_k\End_K(E)\); in the corresponding descent action,
\(H\) fixes \(\Delta_0\) pointwise.
\end{theorem}

Theorem~\ref{thm:parameters} applies to the stated \(U,W,h,M\).
In particular, \(X\) is geometrically integral, its singular locus
has codimension greater than \(h\), and its punctured smooth cone
\(Y^o\) is \((M+2h)\)-connected after a complex embedding.
Multiplication of each block by a \(b\)-th root of unity preserves
the defining equations and linearizes \(\cO_X(1)\).
The Picard factors are geometrically integral, so \(K\) is a field.
The action on \(R\) makes the action on \(K\) faithful.  The
fixed-field theorem therefore gives \(\Gal(K/F)=H\).

Choose a point on the constant curve \(C\), and normalize its
universal line bundles along that point.  The generic section
spaces in the theorem have dimension \(e\) by
Proposition~\ref{prop:binary-curves}, so \(\dim_K E=e_1\).
We use left transport \(g_*=(g^{-1})^*\).
Transport of a normalized universal line differs from the universal
line at the transported parameter by a line from the parameter base.
At the generic point an identification is therefore unique up to a
scalar.  Such a scalar disappears on endomorphisms, so transport
gives an actual semilinear action on \(\End_K(E)\).
Galois descent, applied also to multiplication and the unit, gives
the asserted central simple algebra \(\Delta_1\)
\cite[Tag 0CDR]{Stacks}.  This proves all assertions of
Theorem~\ref{thm:division} except that \(\Delta\) is division.
The algebra asserted here is over the generic field \(F\). The section
bundle \(E\) exists on the non-effectivity open in \(B\); no extension
of its endomorphism algebra as an Azumaya algebra over all \(B\times X\)
is needed.

If a central simple algebra has degree \(d\), the dimension of each
finite right module is divisible by \(d\), as one sees after a
splitting extension.  We call the quotient its \emph{reduced rank}.
To prove that the algebra is division, it suffices to show that
every reduced rank is divisible by \(d\): a nonzero proper right ideal
would have reduced rank strictly between \(0\) and \(d\).
Accordingly, fix a finite right \(\Delta\)-module of reduced rank
\(r\).  To prove \(e_0e_1\mid r\), introduce the auxiliary products
over \(k\)
\[
P=\prod_{m=2}^M\Pic^1(C),\qquad Z=B\times P,
\]
with the corresponding transport action of \(H\).
These degree-one Picard factors serve two purposes. Their free action
will make equivariant sheaves descend to a finite quotient, while a
determinant line coupling \(B\) to \(P\) will turn the module's reduced
rank into an Euler characteristic. They are used only in this proof;
the fields \(K,F\) and the algebra \(\Delta\) remain those in
Theorem~\ref{thm:division}. The following proposition records the two
numerical identities that the index argument will compare.

\begin{proposition}\label{prop:rank-to-euler}
For the data of Theorem~\ref{thm:division} and a right
\(\Delta\)-module of reduced rank \(r\), there is a class \(\beta\)
of \(H\)-equivariant coherent sheaves on \(Z\times X\), with a
commuting right \(\Delta_0\)-action, having the following properties.

After any embedding \(k\hookrightarrow\C\), identify
\(\Delta_0\otimes_k\C\) with \(\operatorname{Mat}_{e_0}(\C)\), choose
a primitive matrix idempotent \(\epsilon\), and put
\(\beta'=\beta\epsilon\).
The class \(\beta'\) is still \(H\)-equivariant and is perfect on
\(Z\times X_{\sm}\).  For every \(x\in X_{\sm}(\C)\),
\begin{equation}\label{eq:descent-fiber-index}
\int_Z\ch\bigl(\beta'|_{Z\times x}\bigr)\td(Z)=\pm r.
\end{equation}
The sign depends only on the curve factors, not on the module or on
\(x\).  Before making the complex base change, one has
\begin{equation}\label{eq:descent-twist-divisibility}
p(l):=\frac{\chi(Z\times X,\beta\otimes\cO_X(l))}{e_0}
       \ \in\ e_0e_1\Z\qquad(l\in\Z).
\end{equation}
After complex base change, \(p(l)\) is also the ordinary Euler
characteristic of \(\beta'\otimes\cO_X(l)\).
\end{proposition}

\subsection{Removing the scalar ambiguity in transport}

Pull the given module to \(K\), with its Galois descent action, and
tensor over \(\End_K(E)\) with its standard left module \(E\).
The resulting vector space \(\cF_\eta\), at the generic point of
\[
T=B\times X\times R,
\]
has a commuting right \(\Delta_0\)-action and dimension \(e_0r\).
Tensoring with \(E\) has divided the dimension by \(e_1\).
Its transport has scalar weight one in each of the universal lines
\(A_m\).  We first extend this twisted transport over \(T\).

For \(g\in H\), simultaneous transport on the curve and the parameter
gives two line-bundle families \(g_*A_m,A_m\) representing the same
Picard parameter.  Their Hom bundle is the pullback of a line from
the base. More explicitly, if \(q_m:T\times C\to T\) is projection for
the \(m\)-th curve, put
\[
 J_g=\bigotimes_{m=2}^M(q_m)_*\mathcal{H}om(g_*A_m,A_m).
\]
Each factor is a line bundle, and its pullback is identified with the
corresponding Hom bundle by evaluation. Evaluation and composition give
\begin{equation}\label{eq:hom-line-transport}
J_g\otimes g_*\cF_\eta\simeq\cF_\eta,\qquad
J_g\otimes g_*J_{g'}\simeq J_{gg'}.
\end{equation}
The second isomorphism has its usual associative composition law.
The first is compatible with it: transport of the original module
is genuine, and the only scalar choices occur in the line bundles
used to define \(E\).  The genuine linearization of \(L\) adds no
further ambiguity.

\begin{lemma}\label{lem:twisted-extension}
The generic module \(\cF_\eta\) extends to a coherent
\(\Delta_0\)-module \(\cF\) on \(T\) with compatible isomorphisms
\(J_g\otimes g_*\cF\simeq\cF\).
\end{lemma}
\begin{proof}
Choose finitely many rational sections spanning \(\cF_\eta\), and
let \(\cF_0\) be the coherent subsheaf of rational sections generated
by them and their multiples by a \(k\)-basis of \(\Delta_0\).
For each \(g\), use \eqref{eq:hom-line-transport} to view
\(J_g\otimes g_*\cF_0\) as a coherent subsheaf of rational sections.
Its image is coherent because \(T\) is noetherian.  Let \(\cF\)
be the sum of these finitely many images.  Corrected transport
permutes the summands by the composition law in
\eqref{eq:hom-line-transport}; hence it preserves \(\cF\).
All maps commute with \(\Delta_0\), and their composition identities
hold inside rational sections.  The generic fiber is unchanged.
\end{proof}

Pull \(\cF\) back along the auxiliary factors \(P\), and let \(V_m\)
be their normalized universal degree-one lines. We will cancel the
scalar weight of \(\cF\) with a line bundle on \(Z=B\times P\).
Define
\begin{equation}\label{eq:determinant-correction}
\cD=\bigotimes_{m=2}^M
\left(
 \bigl(\det R\Gamma_C(A_m\otimes V_m)\bigr)^{-1}
 \otimes\det R\Gamma_C(A_m)
 \otimes\det R\Gamma_C(V_m)
\right).
\end{equation}
Here \(R\Gamma_C\) means derived pushforward along the constant
curve, retaining the Picard parameters. The projection is smooth and
proper, and the universal lines are flat over the parameter base;
their derived pushforwards are perfect and commute with base change
\cite[Tag 07VK]{Stacks}.
Their determinant lines use the determinant-of-cohomology formalism of
Knudsen and Mumford \cite{KnudsenMumford76}; the functorial construction
and its compatibility with pullback are recalled in
\cite[Tags 0FJI and 0FJW]{Stacks}.
Each factor in \(\cD\) is the inverse Deligne pairing of \(A_m,V_m\)
tensored with the fixed line \(\det R\Gamma_C(\cO_C)\)
\cite[Construction~7.2, Equation~(7.2.1)]{Deligne87}.
We keep the full determinant expression when defining transport.

\begin{lemma}\label{lem:determinant-descent}
The sheaf \(\cF\otimes\cD\) on \(Z\times X\times R\) has a
genuine \(H\)-equivariant structure commuting with \(\Delta_0\).
Its derived restriction to \(R=0\) defines an equivariant coherent
class \(\beta\) on \(Z\times X\).
Nonequivariantly this class is the tensor product of \(\cD\) with
a class pulled back from \(B\times X\).
\end{lemma}
\begin{proof}
Write \(g_C\) for the genus of \(C\).  For one pair \(A,V\) of
degrees zero and one,
\[
\chi(A)=1-g_C,\qquad
\chi(A\otimes V)=\chi(V)=2-g_C.
\]
A scalar acting on a line inside a determinant of cohomology has
exponent equal to its Euler characteristic.  Thus the corresponding
factor of \(\cD\) has scalar weights
\[
\bigl(-\chi(A\otimes V)+\chi(A),
      -\chi(A\otimes V)+\chi(V)\bigr)=(-1,0)
\]
in \(A,V\).  The first weight cancels the weight one of \(\cF\),
and the second leaves no dependence on identifications of \(V\).
Local choices of line identifications therefore give the same
transport on \(\cF\otimes\cD\). To check the group law, compare the
two successive transports with the composed identifications of the
universal lines. Composition in
\(J_g\otimes g_*J_h\longrightarrow J_{gh}\) is the associative
composition from \(g_*h_*A\) to \(g_*A\) to \(A\).
Functoriality of determinants and the genuine transport of the original
module identify the two maps for these composed choices. Any scalar
change of the \(A\)-identification contributes opposite weights on
\(\cF\) and \(\cD\), and a change of the \(V\)-identification has
weight zero. Thus no scalar cocycle remains, and the maps glue to
genuine equivariance.

The origin of the representation \(R\) is \(H\)-fixed.  Its
immersion in the product has the finite equivariant Koszul
resolution on the linear coordinates of \(R\), even when \(X\)
is singular.  The alternating coherent Tor class of restriction
therefore defines \(\beta\), preserving both actions.  Since
\(\cF\) does not depend on \(P\), the last assertion follows.
\end{proof}

\subsection{Extracting the reduced rank}

Fix a complex embedding.  All varieties in this paragraph are
base changed along it.  The action of \(H\) is now \(\C\)-linear
and fixes the constant algebra \(\Delta_0\otimes_k\C\) pointwise.
Consequently taking the image under right multiplication by
\(\epsilon\) is exact and preserves equivariance.  It divides the
dimension of every finite-dimensional matrix module by \(e_0\).
The product \(B\times X_{\sm}\times R\) is smooth and connected.
On it, \(\cF\) is perfect, and the alternating rank of a locally
free resolution is constantly \(e_0r\).  Derived restriction to
\(R=0\) preserves this rank.  Thus, for each \(x\in X_{\sm}\),
Lemma~\ref{lem:determinant-descent} gives
\begin{equation}\label{eq:beta-fiber-factorization}
\beta'|_{Z\times x}=\operatorname{pr}_B^*\gamma_x\otimes\cD,
\qquad \gamma_x\in K^0(B),\qquad \rk\gamma_x=r.
\end{equation}
This argument uses derived restriction; no flatness of \(\cF\)
over \(R\) is required.

\begin{lemma}\label{lem:mixed-picard-pairing}
For any class \(\gamma\in K^0(B)\),
\[
\int_Z\ch(\operatorname{pr}_B^*\gamma\otimes\cD)\td(Z)
   =\pm\rk\gamma,
\]
where the sign is independent of \(\gamma\).
\end{lemma}
\begin{proof}
It suffices to calculate the determinant correction for one curve.
Write \(a=c_1(A)\), \(v=c_1(V)\), and let \(q\) project off \(C\).
Grothendieck--Riemann--Roch for this smooth proper projection gives
\[
c_1\!\left(\bigl(\det Rq_*(A\otimes V)\bigr)^{-1}
       \otimes\det Rq_*A\otimes\det Rq_*V\right)
       =-q_*(av).
\]
Indeed the signed exponential sum
\(-e^{a+v}+e^a+e^v\) has degree-two part zero and degree-four
part \(-av\); a curve has no degree-four Todd term.

Normalization at the chosen curve point removes pure
parameter-space degree-two components from both universal classes.
The class \(a\) has no pure curve component because \(\deg A=0\).
The pure curve component of \(v\) cannot contribute to \(q_*(av)\).
The resulting class is therefore purely mixed between
\(H^1(\Pic^0(C),\Z)\) and \(H^1(\Pic^1(C),\Z)\).
Tensoring by \(\cO_C(-c_0)\), for the chosen curve point \(c_0\),
identifies the torsor \(\Pic^1(C)\) with \(\Pic^0(C)\); different
choices change this identification by translation, which acts trivially
on integral first cohomology. The universal classes identify these
integral lattices with the curve's first-cohomology lattice: pullback by
an Abel map gives
the mixed component of the diagonal class on \(C\times C\), and
the Abel map induces an integral \(H^1\)-isomorphism.
The intersection form on \(H^1(C,\Z)\) is unimodular.  Hence the
mixed class has a unimodular matrix \(\Lambda\).
This is the cohomological Poincar\'e pairing; the rank extraction below
is the numerical counterpart of the Fourier--Mukai calculation in
\cite[Theorem~2.2 and Example~2.6]{Mukai81}.

More explicitly, if \(d=2g_C\) is the real dimension of either
Picard factor and \(b_i,p_j\) are integral bases, a mixed class
\(\theta=\sum_{i,j}\Lambda_{ij}b_i\wedge p_j\) satisfies
\[
\int_{\Pic^0(C)\times\Pic^1(C)}\frac{\theta^d}{d!}
       =\pm\det\Lambda=\pm1.
\]
The factors for distinct curves form separate blocks, so the same
unit calculation holds on \(Z\).
Each Picard tangent bundle is trivial, giving \(\td(Z)=1\).
To obtain top degree on \(P\) from \(\exp(c_1(\cD))\), one must
already use top degree on \(B\), because every mixed term has
degree one on each side.  Thus every positive-degree term of
\(\ch(\gamma)\) contributes zero.  Only its rank remains.
\end{proof}

Combining \eqref{eq:beta-fiber-factorization} with
Lemma~\ref{lem:mixed-picard-pairing} proves
\eqref{eq:descent-fiber-index}.  It remains to obtain the arithmetic
divisibility of the twists over \(k\), where \(\Delta_0\) is still
division.  We need an Euler-characteristic fact that also applies
to the singular variety \(X\).

\begin{lemma}\label{lem:torsor-euler}
Let \(Q\) be a projective scheme over a field, and let
\(\delta:T'\to Q\) be a torsor under a constant finite group
of order \(a\).  For every coherent-sheaf class \(\alpha\) on \(Q\),
\[
\chi(T',\delta^*\alpha)=a\,\chi(Q,\alpha).
\]
\end{lemma}
\begin{proof}
The vector bundle \(\delta_*\cO_{T'}\) has rank \(a\).
Put \(d=[\delta_*\cO_{T'}]-a\in K^0(Q)\).
Multiplication by \(d\) lowers the dimension-of-support filtration
on the coherent-sheaf group \(G_0(Q)\).
To see this, the successive quotients of that filtration are
generated by integral closed supports
\cite[Lemma~42.23.2, Tag 02S9]{Stacks}; on an integral support the two bundles
in \(d\) have the same rank and are isomorphic on a dense open
set.  Localization therefore puts their difference in smaller
support dimension.  In particular this operator is nilpotent.

The torsor becomes the disjoint union of \(a\) copies of \(T'\)
after pullback to \(T'\), so \(\delta^*d=0\).
The projection formula gives, as operators on \(G_0(Q)\),
\[
(a+d)d=0.
\]
Over \(\Q\), the operator \(a+d\) is invertible by a finite
nilpotent expansion, hence \(d\) acts by zero.
Euler characteristic takes values in \(\Z\), so it kills torsion
and \(\chi(Q,d\alpha)=0\).
Finite pushforward and the projection formula now give the
claimed equality.
\end{proof}

\begin{proof}[Proof of Proposition~\ref{prop:rank-to-euler}]
The construction and the first identity have been established
above.  By Proposition~\ref{prop:binary-curves}, \(H\) acts freely
on \(P\), hence on \(Z\times X\).  Its quotient is a projective
scheme, and the quotient morphism is a finite étale \(H\)-torsor;
existence of the scheme quotient and torsor follows from
\cite[Tag 07S7]{Stacks}, and a product of the translates of an
ample line bundle gives projectivity of the quotient.
Equivariant coherent sheaves descend along this torsor, together
with their commuting \(\Delta_0\)-actions
\cite[Proposition~35.5.2, Tag 023R]{Stacks}. The same holds for
\(\beta\otimes\cO_X(l)\), since \(\cO_X(1)\) is linearized.

Each cohomology group of a descended sheaf is a finite right
module over the division algebra \(\Delta_0/k\), so its
\(k\)-dimension is divisible by \(e_0^2\).
Lemma~\ref{lem:torsor-euler} multiplies its Euler characteristic
by \(|H|=e_1\) on passing upstairs.  Applying this term by term
to the coherent class gives
\[
\chi(Z\times X,\beta\otimes\cO_X(l))\in e_0^2e_1\Z,
\]
which proves \eqref{eq:descent-twist-divisibility}.
Finally, field extension preserves coherent cohomology dimensions
by flat base change \cite[Tag 02KH]{Stacks}, even for a complex
embedding that does not fix the original constants.  Over \(\C\),
idempotent reduction divides these dimensions by \(e_0\), proving
the last assertion.
\end{proof}

We have obtained an Euler characteristic equal to the reduced
rank up to sign, together with divisibility for every algebraic
twist.  The geometry of \(X\) prevents an immediate comparison of
these two statements.  The next section makes that comparison
using the connectivity of \(Y^o\), and completes the proof of
Theorem~\ref{thm:division}.
